\chapter{Background and related work}
\label{ch:background}

\section{Thyroid nodules, TI-RADS and the reference standard}
\label{sec:bg-thyroid}

Thyroid nodules are common and most of them are benign. Ultrasound is the first examination, and
risk systems such as ACR TI-RADS~\citep{tessler2017} score a nodule from its appearance:
composition, echogenicity, shape (in particular a taller-than-wide shape), margin and echogenic
foci. The score and the size of the nodule decide whether a fine-needle aspiration (FNA) is
recommended. The cytology of the aspirate, and surgery where needed, establish the diagnosis.

FNA is the reference standard but it is invasive and uses clinical resources, and many aspirated
nodules turn out benign. A system that could clear clearly benign nodules with a known error
rate, and pass all doubtful ones to a clinician, would reduce this load without hiding its
mistakes.

\section{Radiomics}
\label{sec:bg-radiomics}

Radiomics describes a region of an image with many hand-crafted, quantitative features: shape
(size, elongation, roundness, outline), first-order statistics of the intensities inside the
region, and texture features computed from matrices that count how grey levels are arranged
(co-occurrence, run lengths, size zones, dependences, neighbourhood differences). PyRadiomics is a
widely used open implementation~\citep{vangriethuysen2017}. Texture features depend on how the
intensities are discretised into grey levels, and shape features depend on the physical scale of
the image; both matter for TN3K (Section~\ref{sec:methods-extraction}).

Radiomics yields many correlated features for a modest number of cases, so a feature selection
step is usual. Minimum-redundancy maximum-relevance (mRMR)~\citep{peng2005} adds features one by one,
preferring those related to the label and unrelated to the features already chosen.

\section{Language models for radiomics and tabular data}
\label{sec:bg-llm}

A row of tabular data can be written as text, ``feature name: value'', and given to a pre-trained
language model. The idea is that the model may link feature names to knowledge acquired in
pre-training, which a classical model that sees only a column of numbers cannot do. Fine-tuning
all weights of a large model is expensive; low-rank adaptation (LoRA)~\citep{hu2022lora} trains
small added matrices instead, and QLoRA~\citep{dettmers2023qlora} does so on a 4-bit copy of the
frozen model, which makes 7B-parameter models trainable on a single mid-range GPU.

\subsection{Ra et al.\ (2025)}

Ra et al.~\citep{ra2025} apply this idea to mammography. They extract 67 PyRadiomics features
(shape, first-order, GLCM, GLSZM) from segmented masses, select 8 with mRMR, write them as
sentences, and fine-tune Llama-2-7B with LoRA (rank 16) and a classification head. On VinDr-Mammo
(235 test images) their model reaches AUC 0.700 against 0.605--0.727 for classical and tabular
baselines; on INbreast, after retraining on half of the data, 0.862. They also report that
replacing the feature names with generic labels lowers performance, and that their model can be
extended to a new feature set without retraining from scratch. This thesis adapts their pipeline
to thyroid ultrasound; a critical reading of their evaluation is given in
Section~\ref{sec:disc-llm}.

\section{Conformal prediction and risk control}
\label{sec:bg-conformal}

Conformal prediction~\citep{vovk2005} turns any model score into prediction sets with a
finite-sample guarantee: the true label lies in the set with probability at least $1 - \alpha$,
provided the calibration data and the new data are exchangeable. It makes no assumption on the
model or the data distribution. With two classes, a set is either one label (a decision), both
labels (``unsure'') or empty, and its guarantee concerns coverage, not the error rate among the
decisions. Class-conditional (Mondrian) versions give the guarantee separately for each class.

Selective classification lets a model abstain and bounds the error rate among the cases it does
decide~\citep{geifman2017}. Learn then Test (LTT)~\citep{angelopoulos2021ltt} generalises this: any
set of risks, written as expected losses, can be controlled by testing candidate parameters as
hypotheses on calibration data, with a guarantee that holds with probability $1 - \delta$. Conformal
risk control~\citep{angelopoulos2024crc} bounds the expected value of a monotone loss, such as the
share of positives that are missed. All of these assume exchangeability. When the test data are
shifted, weighted versions can restore the guarantee under covariate shift~\citep{tibshirani2019}.

\section{Uncertainty from repeated queries}
\label{sec:bg-queries}

For language models, a common uncertainty signal is agreement across repeated or perturbed
queries: self-consistency over sampled answers~\citep{wang2023selfconsistency}, semantic entropy
over answers with the same meaning~\citep{kuhn2023semantic}, and conformal prediction built on
sampling frequencies when the model's probabilities are not accessible~\citep{su2024api}. A recent
scoping review~\citep{ashby2026} finds that such sampling-based signals often outperform those
based on the model's probabilities, mostly on general language benchmarks. Noorani et
al.~\citep{noorani2025} study how many queries are needed for a conformal guarantee. Whether these
findings carry over to a fine-tuned classifier reading radiomics prompts is one of the questions of
this thesis.
